\documentclass[12pt, letterpaper]{article}
\usepackage[utf8]{inputenc}
\usepackage[spanish, es-tabla]{babel}
\usepackage{geometry}
\geometry{top=2.5cm, bottom=2.5cm, left=2.5cm, right=2.5cm}
\usepackage{lmodern}
\usepackage{graphicx}
\graphicspath{{recursos/}{./}}
\usepackage{booktabs}
\usepackage{tabularx}
\usepackage[table]{xcolor}
\usepackage{hyperref}
\usepackage{caption}
\usepackage{float}
\usepackage{listings}
\usepackage[most]{tcolorbox}
\usepackage{tikz}
\usetikzlibrary{positioning,arrows.meta}

% Configuración de marcos para imágenes
\setlength{\fboxsep}{0pt}
\setlength{\fboxrule}{0.6pt}

% Configuración de colores y estilo para código
\definecolor{codebg}{gray}{0.96}
\definecolor{codecomment}{HTML}{2F5D8C}
\definecolor{codestring}{HTML}{C00000}
\definecolor{codekeyword}{HTML}{007F00}

\lstdefinelanguage{mongosh}{
  morekeywords={const,let,var,function,return,true,false,null},
  morecomment=[l]{//},
  morestring=[b]",
  sensitive=true}

\lstdefinestyle{lab}{
  basicstyle=\fontsize{9}{11}\selectfont\ttfamily,
  backgroundcolor=\color{codebg},
  numbers=left, numberstyle=\fontsize{7}{8}\selectfont\color{gray}, numbersep=7pt,
  commentstyle=\itshape\color{codecomment},
  keywordstyle=\color{codekeyword},
  stringstyle=\color{codestring},
  showstringspaces=false, breaklines=true, keepspaces=true,
  frame=single, framerule=0pt, rulecolor=\color{codebg}, framesep=5pt,
  xleftmargin=14pt, xrightmargin=0pt,
  aboveskip=8pt, belowskip=8pt,
  extendedchars=true,
  literate={á}{{\'a}}1 {é}{{\'e}}1 {í}{{\'i}}1 {ó}{{\'o}}1 {ú}{{\'u}}1
           {Á}{{\'A}}1 {É}{{\'E}}1 {Í}{{\'I}}1 {Ó}{{\'O}}1 {Ú}{{\'U}}1
           {ñ}{{\~n}}1 {Ñ}{{\~N}}1 {ü}{{\"u}}1 {¿}{{?`}}1 {¡}{{!`}}1
}

% ----- columnas de tabla
\newcolumntype{Y}[1]{>{\raggedright\arraybackslash}p{\dimexpr #1\linewidth-2\tabcolsep\relax}}

\hypersetup{
    colorlinks=true,
    linkcolor=blue!70!black,
    citecolor=blue!70!black,
    urlcolor=blue!70!black,
    pdftitle={Proyecto Final - Entrega 2: Fronteras PDET e integración espacial},
    pdfauthor={Samuel Tovar, Juan Romero, Carlos Pinzón Caicedo, Alejandro Padilla Gaviria}
}

\begin{document}

% ==========================================
% PORTADA ACADÉMICA
% ==========================================
\begin{titlepage}
    \centering
    \vspace*{0.5cm}

    {\huge \textbf{Administración de Bases de Datos}}\\[0.4cm]
    {\Large \textbf{Proyecto Final: Potencial Solar de Techos en Municipios PDET}}\\[1.2cm]

    \begin{figure}[H]
        \centering
        \includegraphics[width=0.5\linewidth]{logo_javeriana.png}
        \label{fig:logo}
    \end{figure}

    {\large \textbf{Estudiantes:}}\\[0.2cm]
    {\large Samuel Tovar}\\[0.1cm]
    {\large Juan Romero}\\[0.1cm]
    {\large Carlos Pinzón Caicedo}\\[0.1cm]
    {\large Alejandro Padilla Gaviria}\\[0.8cm]

    {\large \textbf{Entrega 2:}}\\[0.2cm]
    {\large Fronteras PDET (DANE MGN 2025) e Integración Espacial en MongoDB}\\[1.2cm]

    \vfill

    {\Large \textbf{Pontificia Universidad Javeriana}}\\[0.6cm]

    {\large 9 de noviembre de 2026}
\end{titlepage}

\clearpage

% ==========================================
% ÍNDICES
% ==========================================
\tableofcontents
\vspace{0.8cm}
\listoffigures
\vspace{0.8cm}
\listoftables
\clearpage

% ==========================================
% SECCIÓN 1: INTRODUCCIÓN
% ==========================================
\section{Introducción}

Esta entrega integra al proyecto los límites de los 170 municipios PDET, tomados de la capa municipal del Marco Geoestadístico Nacional 2025 del DANE, y los deja cargados en MongoDB con índices espaciales. Es la base sobre la que, en las entregas siguientes, se asignarán las huellas de edificios a cada municipio y se calcularán el número de techos y su área. El enunciado evalúa tres aspectos, que el Cuadro~\ref{tab:criterios} relaciona con las secciones del informe y con los archivos del repositorio (\url{https://github.com/SamT81/Proyecto-ABD}).

\begin{table}[H]
\centering
\small
\rowcolors{2}{gray!12}{white}
\begin{tabularx}{\textwidth}{Y{0.25} Y{0.14} Y{0.38} Y{0.23}}
\toprule
\textbf{Criterio del enunciado} & \textbf{Sección} & \textbf{Archivos principales} & \textbf{Rama} \\
\midrule
Adquisición y verificación de los datos & 2 & \texttt{00a\_download\_mgn.py}, \texttt{00b\_build\_pdet\_catalog.py}, \texttt{00c\_filter\_pdet.py}, \texttt{docs/calidad/} & \texttt{feat/pdet-acquisition} \\
Integridad y formato para la ingesta & 3 & \texttt{01\_load\_municipalities.py}, \texttt{utils/geo\_integrity.py} & \texttt{feat/pdet-integrity-etl} \\
Integración espacial NoSQL (índices) & 4 & \texttt{load\_to\_mongo.py}, \texttt{pdet\_mongo.py}, \texttt{spatial\_evidence.py} & \texttt{feat/spatial-mongo} \\
\bottomrule
\end{tabularx}
\caption{Criterios de evaluación de la Entrega 2 y su ubicación en el informe y en el repositorio.}
\label{tab:criterios}
\end{table}

\subsection{Flujo de la entrega}
La Figura~\ref{fig:flujo} muestra el recorrido de los datos. Cada etapa verifica su entrada y se detiene con un error si un control falla, de modo que no llegan datos dudosos a la etapa siguiente.

\begin{figure}[H]
\centering
\begin{tikzpicture}[
  node distance=5mm and 7mm,
  box/.style={draw, rounded corners=2pt, fill=blue!6, align=center, font=\footnotesize,
              minimum width=2.6cm, minimum height=1.05cm},
  arr/.style={-{Stealth[length=2mm]}, thick}]
  \node[box] (a) {Descarga y\\verificación\\\texttt{00a}};
  \node[box, right=of a] (b) {Filtro PDET y\\controles\\\texttt{00c}};
  \node[box, right=of b] (c) {ETL de integridad\\\texttt{01\_load\_}\\\texttt{municipalities}};
  \node[box, right=of c] (d) {Carga e índices\\\texttt{load\_to\_}\\\texttt{mongo}};
  \node[box, below=of d, fill=green!8] (e) {Evidencia\\\texttt{spatial\_}\\\texttt{evidence}};
  \node[font=\scriptsize, below=1mm of a, align=center] {MGN2025\\(DANE)};
  \node[font=\scriptsize, below=1mm of b, align=center] {\texttt{.gpkg}\\170 municipios};
  \node[font=\scriptsize, below=1mm of c, align=center] {\texttt{.geojson}\\EPSG:4326};
  \draw[arr] (a)--(b); \draw[arr] (b)--(c); \draw[arr] (c)--(d); \draw[arr] (d)--(e);
\end{tikzpicture}
\caption{Flujo de la Entrega 2: de la descarga del MGN2025 a la evidencia de los índices en MongoDB.}
\label{fig:flujo}
\end{figure}

% ==========================================
% SECCIÓN 2: ADQUISICIÓN Y VERIFICACIÓN
% ==========================================
\newpage

\section{Adquisición y Verificación de los Datos}

\subsection{Fuentes y mecanismo de descarga}
Los límites municipales se toman del Marco Geoestadístico Nacional 2025 del DANE. En el plan de la Entrega~1 se previó descargar el ZIP completo (1.5~GB, todos los niveles geográficos), pero el geoportal publica el nivel municipio por separado (\texttt{MGN2025\_MPIO\_GRAFICO.zip}, 71.6~MB, Shapefile en WGS84) con la misma versión y la misma capa, por lo que se descarga solo ese archivo. El listado de los 170 municipios PDET y su subregión se toma de la Agencia de Renovación del Territorio (Decreto 893 de 2017).

La descarga la hace \texttt{src/00a\_download\_mgn.py}: lee el archivo por bloques de 1~MB a un temporal \texttt{.part}, reintenta y reanuda con cabecera \texttt{Range} si la conexión se corta, y solo renombra el archivo cuando su tamaño coincide con el \texttt{Content-Length} del servidor. La URL, las rutas y los valores de referencia están en \texttt{config/acquisition.yaml}, no en el código.

\begin{lstlisting}[style=lab]
python src/00a_download_mgn.py      # descarga y verifica (si ya existe, solo verifica)
python src/00c_filter_pdet.py       # filtro PDET y controles -> data/mgn_pdet_2025.gpkg
python src/01_load_municipalities.py
\end{lstlisting}

\subsection{Verificación de integridad}
El DANE no publica un \emph{checksum} oficial. Por eso, en la primera descarga (08/11/2026) se fijó en la configuración el SHA-256 del ZIP (\texttt{a3c01393\ldots bc141f}), su tamaño (71\,576\,494 bytes) y las cabeceras \texttt{Last-Modified} y \texttt{ETag}. Cada corrida recalcula el SHA-256 y falla si no coincide, lo que garantiza que todo el equipo procese exactamente los mismos bytes y detecta si el DANE reemplaza el archivo. Además se valida el CRC de los ocho archivos del ZIP y la presencia de los componentes del Shapefile (\texttt{.shp}, \texttt{.shx}, \texttt{.dbf}, \texttt{.prj}, \texttt{.cpg}). El SHA-256 de cada archivo extraído queda en \texttt{docs/calidad/manifest\_mgn.json} y se vuelve a comprobar antes del filtrado.

\subsection{Filtro PDET y controles de calidad}
\texttt{src/00c\_filter\_pdet.py} cruza la capa municipal (1\,122 unidades) con el listado PDET por código DIVIPOLA. En el Excel de la ART los códigos de Antioquia aparecen sin el cero inicial (\texttt{5490} en lugar de \texttt{05490}), así que se normalizan a cinco dígitos antes del cruce. El Cuadro~\ref{tab:qc_mgn} resume los 32 controles; el detalle está en \texttt{docs/calidad/qc\_mgn\_pdet.md}.

\begin{table}[H]
\centering
\small
\rowcolors{2}{gray!12}{white}
\begin{tabularx}{\textwidth}{l X c}
\toprule
\textbf{Grupo} & \textbf{Controles} & \textbf{Resultado} \\
\midrule
Fuente & SHA-256 del ZIP verificado; archivos extraídos sin cambios & 2 OK \\
Capa MGN & Campos requeridos; CRS EPSG:4326; 1\,122 unidades; código único, de 5 dígitos y coherente con el departamento; sin nulos; extensión dentro de Colombia & 8 OK \\
Listado PDET & 170 municipios; códigos únicos y de 5 dígitos; 16 subregiones; prefijo coherente; sin nulos & 6 OK \\
Cruce & 170/170 códigos encontrados; departamento coincide; nombre coincide & 2 OK, 1 AVISO \\
Geometría & 170 polígonos únicos, no vacíos y válidos; área EPSG:9377 vs.\ \texttt{mpio\_narea} (diferencia 0.000\,\%); sin solapes (339 pares vecinos); geometría más pesada 2.2~MB (límite BSON 16~MB) & 10 OK \\
Salida & GeoPackage releído: 170 registros, EPSG:4326, sin nulos & 3 OK \\
\bottomrule
\end{tabularx}
\caption{Controles de calidad sobre el MGN2025 filtrado a los municipios PDET (31 OK, 1 AVISO, 0 FALLA).}
\label{tab:qc_mgn}
\end{table}

El único aviso corresponde a seis municipios cuyo nombre en el listado de la ART difiere del nombre oficial del MGN2025. Como el cruce es por código, no afecta el resultado; se conserva el nombre del MGN. Los 170 municipios suman 389\,182~km$^2$ en 16 subregiones. La salida, \texttt{data/mgn\_pdet\_2025.gpkg} (EPSG:4326, geometrías \texttt{MultiPolygon}), es la entrada del ETL de integridad.

% ==========================================
% SECCIÓN 3: INTEGRIDAD GEOMÉTRICA Y ETL
% ==========================================
\section{Integridad Geométrica y Formateo del ETL}

\subsection{Normalización espacial y reparación topológica}
Una vez filtrada la capa municipal a los 170 municipios PDET (\texttt{data/mgn\_pdet\_2025.gpkg}), el módulo \texttt{src/utils/geo\_integrity.py} y el orquestador \texttt{src/01\_load\_municipalities.py} aseguran que cada registro cumpla los requisitos estrictos del índice \texttt{2dsphere} de MongoDB y del diccionario de datos definido en la Entrega~1:

\begin{itemize}
  \item \textbf{Reparación topológica y tipo uniforme:} toda geometría se valida y repara con \texttt{shapely.validation.make\_valid}, descartando componentes degenerados y promoviendo polígonos simples a \texttt{MultiPolygon}. Asimismo, se orientan los anillos exteriores en sentido antihorario y los huecos en sentido horario (regla de la mano derecha) para evitar errores de indexación en coordenadas esféricas.
  \item \textbf{Doble sistema de referencia:} el área oficial en kilómetros cuadrados (\texttt{area\_km2}) se calcula proyectando temporalmente a \textbf{EPSG:9377} (MAGNA-SIRGAS Origen Nacional) antes de transformar definitivamente las geometrías a \textbf{EPSG:4326} (WGS84), redondeando coordenadas a 7 decimales ($\approx 1.1$~cm de precisión) para reducir el tamaño BSON sin pérdida geométrica apreciable y evitando auto-intersecciones de anillos.
  \newpage
  \item \textbf{Validación de dominio y tamaño BSON:} se verifica que los 170 códigos DIVIPOLA tengan 5 dígitos únicos, sin campos obligatorios nulos, que sus envolventes caigan dentro de la caja geográfica de Colombia ($[-82^\circ, -5^\circ, -66^\circ, 14^\circ]$) y que ningún documento supere el umbral preventivo de 15.5~MB frente al límite de 16~MB de MongoDB.
\end{itemize}

\subsection{Salida estandarizada e ingesta idempotente}
El script \texttt{src/01\_load\_municipalities.py} genera el archivo limpio \texttt{data/processed/pdet\_municipios.geojson} (consumido por la integración espacial), guarda la bitácora de validación en \texttt{docs/calidad/etl\_integrity\_report.json} y permite la carga directa e idempotente en \texttt{municipalities} mediante operaciones \texttt{ReplaceOne(..., upsert=True)} indexadas por \texttt{\_id} (DIVIPOLA).

El ETL se validó con los datos reales del MGN2025: leyó 170 municipios, los 170 resultaron válidos (0 rechazados, 0 geometrías descartadas), en 16 subregiones y con un área total de 389\,182.04~km$^2$ calculada en EPSG:9377. El documento más grande pesa unos 1.24~MB, muy por debajo del límite BSON de 16~MB. Su bitácora queda en \texttt{docs/calidad/etl\_integrity\_report.json}, con el SHA-256 del archivo de entrada y del GeoJSON de salida.

% ==========================================
% SECCIÓN 4: INTEGRACIÓN ESPACIAL EN MONGODB
% ==========================================
\section{Integración Espacial en MongoDB}

\subsection{Carga e índices}
Los municipios PDET limpios (GeoJSON en EPSG:4326, producido por el ETL de integridad) se cargan en la colección \texttt{municipalities} con \texttt{src/load\_to\_mongo.py}. La carga es idempotente (\emph{upsert} por \texttt{\_id}, el código DIVIPOLA), las geometrías se normalizan a \texttt{MultiPolygon} con la regla de la mano derecha y se rechazan las inválidas o ubicadas fuera de la caja de Colombia. Los índices se crean antes de insertar, de modo que MongoDB también rechaza cualquier geometría que no pueda indexar.

\begin{lstlisting}[style=lab,language=mongosh]
db.municipalities.createIndex({ geometry: "2dsphere" }, { name: "geometry_2dsphere" })
db.municipalities.createIndex({ pdet_subregion: 1 },    { name: "pdet_subregion_1" })
\end{lstlisting}

La carga queda registrada en \texttt{ingestion\_log} (archivo, filas leídas, insertadas y rechazadas, SHA-256 del archivo y fecha). Resultado: 170 municipios cargados de 170 esperados, 16 subregiones, 0 rechazados.

\subsection{Evidencia del uso de los índices}
Se ejecutaron consultas \texttt{\$geoIntersects}, \texttt{\$geoWithin} y \texttt{\$near} con \texttt{explain("executionStats")}, comparando el plan con el índice \texttt{2dsphere} frente a un recorrido completo forzado con \texttt{hint(\{\$natural: 1\})}. El Cuadro~\ref{tab:evidencia} resume los planes ganadores y los documentos examinados.

\begin{table}[H]
\centering
\small
\rowcolors{2}{gray!12}{white}
\begin{tabularx}{\textwidth}{X l r r r}
\toprule
\textbf{Consulta} & \textbf{Plan ganador} & \textbf{Docs (índice)} & \textbf{Docs (sin índice)} & \textbf{Devueltos} \\
\midrule
\texttt{punto\_en\_poligono[05031]} & FETCH/IXSCAN & 1 & 170 & 1 \\
\texttt{punto\_en\_poligono[13688]} & FETCH/IXSCAN & 1 & 170 & 1 \\
\texttt{punto\_en\_poligono[19548]} & FETCH/IXSCAN & 2 & 170 & 1 \\
\texttt{punto\_en\_poligono[47001]} & FETCH/IXSCAN & 1 & 170 & 1 \\
\texttt{punto\_en\_poligono[54800]} & FETCH/IXSCAN & 1 & 170 & 1 \\
\texttt{punto\_en\_el\_oceano} & FETCH/IXSCAN & 0 & 170 & 0 \\
\texttt{ventana\_geoWithin} & FETCH/IXSCAN & 34 & 170 & 18 \\
\texttt{vecinos\_geoIntersects} & FETCH/IXSCAN & 5 & 170 & 1 \\
\texttt{cercanos\_near} & GEO\_NEAR\_2DSPHERE... & 335 & -- & 18 \\
\bottomrule
\end{tabularx}
\caption{Planes de ejecución de las consultas espaciales sobre \texttt{municipalities}.}
\label{tab:evidencia}
\end{table}

Los controles automáticos (índice presente, número de municipios y subregiones, códigos DIVIPOLA de cinco dígitos, tipo \texttt{MultiPolygon}, y consultas punto-en-polígono cuyo resultado se verifica contra el municipio de origen) se reportan en \texttt{docs/evidencia/spatial\_evidence.md}. Con 170 documentos la diferencia de tiempo es pequeña, pero el plan cambia de \texttt{COLLSCAN} a \texttt{IXSCAN}; la ventaja crecerá exponencialmente en las colecciones de edificios (Entrega 3), con millones de documentos.

% ==========================================
% SECCIÓN 5: CONCLUSIONES Y SIGUIENTES PASOS
% ==========================================
\section{Conclusiones y Siguientes Pasos}

\begin{itemize}
  \item Los límites de los 170 municipios PDET provienen del MGN2025 con una descarga reproducible y verificada por SHA-256, y pasaron 31 controles de calidad sin fallas (un aviso por diferencias de nombre con el listado de la ART, que no afecta el cruce por código DIVIPOLA).
  \item El ETL entrega geometrías válidas, orientadas, en EPSG:4326, con el área calculada en EPSG:9377, lista para el índice \texttt{2dsphere}; no fue necesario simplificar ni partir ningún municipio tras ajustar a 7 decimales.
  \item La integración espacial crea los índices antes de la carga, registra cada ingesta en \texttt{ingestion\_log} y automatiza la evidencia con \texttt{explain()} para comparar el plan con índice frente al recorrido completo.
  \item \textbf{Limitación:} el tamaño de esta colección (170 documentos) hace que las diferencias de tiempo entre \texttt{IXSCAN} y \texttt{COLLSCAN} sean pequeñas; el beneficio del índice se medirá con las colecciones de edificios de la Entrega~3.
  \item \textbf{Siguientes pasos:} cargar las huellas de Microsoft y Google por lotes, asignar a cada edificio su \texttt{municipality\_code} con el punto representativo, y hacer el análisis exploratorio (Entrega~3).
\end{itemize}

\end{document}
